\documentclass[journal,twoside]{IEEEtran}

\usepackage[T1]{fontenc}
\usepackage{cite}
\usepackage{amsmath,amssymb}
\usepackage{booktabs}
\usepackage{graphicx}
\usepackage{hyperref}
\usepackage{microtype}

% All flash footprint values in this paper are verified measurements
% from: arm-none-eabi-gcc -mcpu=cortex-m4 -mthumb -std=c99 -O2 -c <file>.c
%       arm-none-eabi-size <file>.o  (reading the .text column)

\begin{document}

\title{ECG Arrhythmia Classification for Microcontroller Deployment: \\
       An Inter-Patient Pipeline with Verified C~Export and ARM~Flash~Profiling}

\author{Balajee,~\IEEEmembership{Student,}~Manipal University Jaipur%
\thanks{Manuscript received \today. Code and model artifacts: 
\url{https://github.com/bl4zeh3x/tinyml-explorations}.}}

\markboth{}{ECG Arrhythmia Classification for Microcontroller Deployment}

\maketitle

% ── Abstract ──────────────────────────────────────────────────────────────────
\begin{abstract}
We present a complete, reproducible pipeline for training ECG arrhythmia
classifiers deployable to ARM Cortex-M microcontrollers without external
inference libraries. Two methodological commitments distinguish this work
from the majority of open implementations. First, we follow the
standard inter-patient evaluation protocol of de~Chazal~et~al.\ (2004),
training exclusively on patients in dataset DS1 and evaluating exactly once
on the held-out patients of DS2; intra-patient cross-validation, which
dominates publicly available code, inflates reported accuracy through
patient-specific morphology leakage. Second, we verify that the deployed
artifact is correct: trained models are exported to dependency-free C via
\texttt{m2cgen}, cross-compiled with \texttt{arm-none-eabi-gcc}, and
confirmed to produce bit-identical predictions to their Python counterparts
on matched input rows before any hardware deployment is attempted. A
24-feature representation combining morphological waveform statistics,
frequency-domain band energies, and RR-interval timing is extracted per
beat; the RR features are included specifically to address the known
morphology-only limitation for supraventricular ectopic (S) detection.
Three model tiers are trained: a 200-tree random forest (\emph{Reference}),
a 10-tree random forest with bounded depth (\emph{Edge}), and a single
depth-6 decision tree (\emph{Tiny}). On DS2, the Reference model achieves
0.920 accuracy and 0.518 reliable macro-F1 (four-class, excluding the
near-empty unclassifiable class); the Tiny model achieves 0.708 accuracy
in 3{,}656 bytes of ARM Cortex-M4 compiled code, while the Edge model
achieves 0.830 accuracy in 94{,}252 bytes. All code,
tests, trained model artifacts, and generated C files are publicly
available.
\end{abstract}

\begin{IEEEkeywords}
ECG arrhythmia, TinyML, edge inference, ARM Cortex-M, inter-patient
evaluation, automatic beat classification, AAMI EC57, MIT-BIH
\end{IEEEkeywords}

% ── I. Introduction ───────────────────────────────────────────────────────────
\section{Introduction}
\IEEEPARstart{C}{ardiac} arrhythmias affect an estimated 1.5--5\% of the
global population~\cite{ZoniBerisso2014}. Continuous long-term monitoring with
wearable Holter monitors or cardiac patches provides substantially richer
diagnostic information than a resting ECG, but demands low-power
autonomy: a microcontroller-class processor performing local inference
rather than streaming raw data to a cloud backend. This constraint
motivates the TinyML framing~\cite{Warden2019}: train the classifier on
a workstation, then deploy only the inference path to the resource-limited
device.

Two persistent problems limit the practical value of most published ECG
classifiers. First, evaluation is frequently performed with intra-patient
cross-validation, in which different beats from the \emph{same} patient
may appear in both training and test folds. The classifier can then
exploit patient-specific baseline QRS morphology rather than learning
class-general arrhythmia signatures, producing inflated accuracy figures
that do not transfer to unseen patients~\cite{deChazal2004}.
Second, the gap between a trained model and a deployed model is typically
unverified: code exported to a target language may diverge from the
original due to floating-point ordering differences, language-specific
rounding, or silent API mismatches.

This paper addresses both problems concurrently. We implement the
standard inter-patient evaluation protocol from de~Chazal~et~al.\
\cite{deChazal2004} exactly, using their published DS1/DS2 patient
partition of the MIT-BIH Arrhythmia Database. We then export trained
models to self-contained C code, cross-compile for ARM Cortex-M4, and
verify prediction parity between Python and C before reporting any result.
The pipeline is fully reproducible: every number in this paper can be
regenerated from the public repository by running a single entry-point
script.

\subsection*{Contributions}
\begin{itemize}
  \item A complete, single-entry-point training and evaluation pipeline
        following the de~Chazal inter-patient protocol, including RR-interval
        timing features not present in the original work, with verified
        per-class precision, recall, and F1 for all five AAMI classes.
  \item Automated export of trained sklearn classifiers to dependency-free
        C inference code via \texttt{m2cgen}, with a pytest-based C/Python
        prediction-parity test suite that compiles and runs the C binary
        against matched input rows.
  \item Measured ARM Cortex-M4 flash footprints (via \texttt{arm-none-eabi-size})
        for three model tiers positioned across the size--accuracy tradeoff,
        demonstrating deployability to commodity embedded targets.
\end{itemize}

% ── II. Background ────────────────────────────────────────────────────────────
\section{Background and Related Work}

\subsection{The MIT-BIH Arrhythmia Database}
The MIT-BIH Arrhythmia Database~\cite{Moody2001,Goldberger2000} contains
48 half-hour ambulatory ECG recordings sampled at 360~Hz. Each recording
is annotated beat-by-beat by two cardiologists. Following the AAMI EC57
standard~\cite{AAMI2008}, beat annotations are mapped to five classes:
Normal (N), Supraventricular ectopic (S), Ventricular ectopic (V),
Fusion (F), and Unclassifiable (Q). The unclassifiable class is drawn
almost entirely from paced-beat patients (records 102, 104, 107, 217),
which are excluded from both training and test sets in the standard
inter-patient protocol, leaving near-zero representation of Q under this
evaluation regime.

\subsection{Inter-Patient Evaluation}
The pivotal methodological reference is de~Chazal~et~al.~\cite{deChazal2004},
who defined a patient-disjoint partition of MIT-BIH: DS1 (22 records for
training) and DS2 (22 records for test). This partition is now used as the
standard comparison point in the AAMI beat classification literature.
Prior work using intra-patient evaluation, including several high-citation
CNN and LSTM papers~\cite{Kiranyaz2016,Acharya2017}, reports
substantially higher accuracy because patient-specific QRS morphology
partially identifies individuals rather than arrhythmia classes.

\subsection{TinyML and Edge Inference}
Deploying inference to MCUs without OS or ML runtime support has been
addressed by frameworks including TensorFlow Lite for Microcontrollers
(TFLM)~\cite{David2021} and Edge Impulse. These target neural networks
and carry a runtime overhead measured in kilobytes. For classical
tree-based models, code generators such as \texttt{m2cgen}~\cite{m2cgen}
emit plain C with no dependencies beyond the standard math library.
The trade-off — smaller and more predictable code at some accuracy cost
— is the central design axis of this work.

% ── III. Methods ──────────────────────────────────────────────────────────────
\section{Methods}

\subsection{Dataset and Evaluation Protocol}
We use the MIT-BIH Arrhythmia Database, downloaded programmatically via
the \texttt{wfdb} Python package~\cite{wfdb}. We implement the de~Chazal
inter-patient split exactly:

\medskip
\noindent\textbf{DS1 (train):} 101, 106, 108, 109, 112, 114, 115, 116,
118, 119, 122, 124, 201, 203, 205, 207, 208, 209, 215, 220, 223, 230.

\noindent\textbf{DS2 (test):} 100, 103, 105, 111, 113, 117, 121, 123,
200, 202, 210, 212, 213, 214, 219, 221, 222, 228, 231, 232, 233, 234.
\medskip

Models are trained exclusively on DS1 and evaluated exactly once on DS2.
No hyperparameter search is conducted; model configurations are fixed in
advance based on the deployment size targets described below.

\subsection{Beat Segmentation}
Each beat is represented as a 360-sample window (1~second at 360~Hz)
centred on the annotated R-peak. Beats with windows extending outside the
recording boundary are discarded. Each window is independently z-score
normalised (zero mean, unit variance) so that absolute sensor gain
differences across leads do not contribute to the classification.

\subsection{Feature Extraction}
We extract 24 scalar features per beat, arranged in four groups:

\subsubsection{Amplitude features (8)}
R-peak amplitude; global maximum and minimum; peak-to-peak range;
Q-wave minimum (40 samples left of R-peak); S-wave minimum (40 samples
right of R-peak); R$-$Q and R$-$S differences.

\subsubsection{Statistical features (6)}
Mean, standard deviation, total signal energy, skewness (non-excess),
kurtosis (non-excess), zero-crossing rate.

\subsubsection{Frequency-domain features (6)}
FFT band powers in four clinically relevant bands (0--5, 5--15, 15--40,
40--100~Hz), normalised by total spectral power; dominant frequency
(first positive-frequency bin excluding DC); spectral centroid.

\subsubsection{RR-interval timing features (4)}
Pre-RR interval (time from previous annotated R-peak to current);
post-RR interval (time to next annotated R-peak); pre/post ratio;
rolling local mean of the previous five pre-RR intervals.
These features were added after observing that morphology-only
classification produces recall below 10\% for the S class:
supraventricular ectopic beats are frequently near-identical to Normal
beats in waveform shape, but are defined clinically by premature
firing — a timing, not a morphology, phenomenon. RR features are computed
per-record and reset at record boundaries so that timing information never
bridges across patients.

The 24 features reduce each 1440-byte double-precision waveform window
(360 samples $\times$ 4 bytes) to a 96-byte feature vector (24 floats),
matching the compute budget available for real-time extraction on
a Cortex-M4 at typical ECG inference cadences.

\subsection{Model Tiers}

Three scikit-learn~\cite{scikit-learn} classifiers are trained with
\texttt{class\_weight=`balanced'} to address the severe class imbalance
(approximately 75\% Normal beats in both DS1 and DS2):

\begin{itemize}
  \item \textbf{Reference}: \texttt{RandomForestClassifier(n\_estimators=200)}
        — maximum-accuracy baseline; not intended for embedded deployment.
  \item \textbf{Edge}: \texttt{RandomForestClassifier(n\_estimators=10,
        max\_depth=8)} — targeting Cortex-M4 class devices with $\ge$128~KB
        flash (e.g.\ STM32L4 series with 256~KB flash); measured footprint
        94{,}252~bytes.
  \item \textbf{Tiny}: \texttt{DecisionTreeClassifier(max\_depth=6)}
        — targeting severely resource-constrained devices such as the
        ATmega328P (Arduino Uno, 32~KB flash); measured footprint
        3{,}656~bytes.
\end{itemize}

No cross-validation is used for model selection. DS2 is touched exactly
once, after all training decisions are finalised.

\subsection{C Export and Deployment Verification}

Each model is exported to C using \texttt{m2cgen}~\cite{m2cgen}, which
walks the fitted estimator and emits a self-contained \texttt{score()}
function:

\begin{verbatim}
void score(double *input, double *output);
\end{verbatim}

\noindent The function accepts a 24-element feature vector and populates a
5-element class probability array. No dynamic allocation, no external
libraries, no floating-point library beyond what the target FPU natively
supports. The Edge and Tiny models are exported; the Reference model is
not, as it is not a deployment target.

\textbf{Parity verification.} A pytest suite compiles each generated C
file against a C harness using the system \texttt{gcc} toolchain and
compares the argmax prediction for 10 real feature rows against the
corresponding \texttt{sklearn} prediction. All parity tests pass for both
exported models, confirming that the C inference and the Python model are
numerically equivalent on these inputs.

\textbf{ARM compilation.} The verified C files are then cross-compiled for
bare-metal ARM Cortex-M4:

\begin{verbatim}
arm-none-eabi-gcc -mcpu=cortex-m4 \
  -mthumb -std=c99 -O2 -c ecg_<model>.c
arm-none-eabi-size ecg_<model>.o
\end{verbatim}

\noindent Flash footprints are read from the \texttt{.text} segment of the
resulting object file.

\subsection{Metrics}

We report DS2 accuracy, strict macro-F1 (all five AAMI classes), and
\emph{reliable} macro-F1, which excludes classes with fewer than 30
training examples from the average. Only the Q class falls below this
threshold (8 training examples in DS1, drawn almost entirely from
paced-beat patients that the standard protocol excludes). The excluded set
and its threshold are always reported alongside the reliable score; this
is not a post-hoc selection but a pre-registered decision motivated by the
statistical unreliability of class-level metrics computed from
single-digit training support. Full per-class precision, recall, and F1
are reported for all three models.

% ── IV. Results ───────────────────────────────────────────────────────────────
\section{Results}

\subsection{Overall Performance}

Table~\ref{tab:overall} summarises DS2 held-out performance for all
three model tiers together with ARM Cortex-M4 flash footprints.

\begin{table}[htbp]
\caption{DS2 Held-Out Performance and ARM Cortex-M4 Flash Footprints}
\label{tab:overall}
\centering
\begin{tabular}{lccccc}
\toprule
Model & Acc. & F1 (strict) & F1 (reliable) & Nodes & .text \\
\midrule
Reference & 0.920 & 0.415 & 0.518 & 283{,}328 & --- \\
Edge      & 0.830 & 0.409 & 0.511 & $\sim$3{,}296 & 94{,}252 \\
Tiny      & 0.708 & 0.324 & 0.405 & 121       & 3{,}656 \\
\bottomrule
\end{tabular}
\vspace{2pt}
\small{Reliable F1 excludes Unclassifiable (Q; 8 DS1 training samples).
Flash footprints measured via \texttt{arm-none-eabi-gcc -O2 -mcpu=cortex-m4 -mthumb}.}
\end{table}

\subsection{Per-Class Analysis}

Table~\ref{tab:perclass} shows precision, recall, and F1 per AAMI class
for the Reference model, alongside DS1 and DS2 sample counts.

\begin{table}[htbp]
\caption{Reference Model Per-Class Metrics on DS2}
\label{tab:perclass}
\centering
\begin{tabular}{lrrrrrrr}
\toprule
Class & \multicolumn{2}{c}{Support} & P & R & F1 \\
      & DS1 & DS2 & & & \\
\midrule
Normal (N)         & 45{,}824 & 44{,}218 & 0.957 & 0.958 & 0.958 \\
Supravent. (S)     &    943   &  1{,}836 & 0.254 & 0.094 & 0.138 \\
Ventricular (V)    &  3{,}788 &  3{,}219 & 0.842 & 0.968 & 0.901 \\
Fusion (F)         &    414   &    388   & 0.053 & 0.142 & 0.078 \\
Unclassif. (Q)$^*$&      8   &      7   & 0.000 & 0.000 & 0.000 \\
\bottomrule
\multicolumn{7}{l}{\small $^*$Excluded from reliable macro-F1.}
\end{tabular}
\end{table}

\subsection{Model-Tier Error Profiles}

A notable qualitative difference emerges between model tiers on the
Fusion (F) class: the Edge model achieves 85.1\% recall on F at 9.1\%
precision, while the Reference model achieves 14.2\% recall at 5.3\%
precision. The constrained depth of the Edge model, combined with
\texttt{class\_weight=`balanced'} amplifying F's 414-sample presence,
produces a broad decision region that captures most true F beats at high
false-positive cost. This is not a tuning artefact but a structural
consequence of depth-limited tree ensembles on heavily imbalanced classes:
shallower models spend more of their limited splits on the imbalance
correction, while the Reference model's unconstrained depth allows tighter
per-class boundaries. The practical implication is that the deployment tier
choice also implicitly selects an error profile, not merely an accuracy level.

% ── V. Discussion ─────────────────────────────────────────────────────────────
\section{Discussion}

\subsection{The S-Class Gap}

Supraventricular ectopic beats achieve F1 of 0.138 on the Reference model
despite adequate DS2 support (1,836 test samples). This result is
literature-consistent: de~Chazal~et~al.\ \cite{deChazal2004} reported
positive predictivity around 31\% for SVEB using considerably more
sophisticated wave-boundary features. The fundamental difficulty is
morphological: S-class beats are frequently near-identical to Normal beats
in raw waveform shape; what clinically defines them is prematurity, a
timing phenomenon our RR features partially address but do not resolve.
RR-interval addition improved macro-F1 from 0.316 (morphology only) to
0.415 (with RR features), confirming the features help, while leaving
substantial room for future work.

\subsection{Why Classical ML for This Use Case}

The MCU constraint provides a principled reason to prefer classical
tree-based models over convolutional or recurrent networks. A depth-6
decision tree performs inference as a sequence of scalar comparisons
against learned thresholds — an operation that maps to a handful of
ARM Thumb-2 instructions with no memory bandwidth requirement beyond
the feature vector. CNN-based ECG classifiers achieve stronger absolute
performance but require kilobytes of parameter storage, floating-point
convolution, and in most published implementations a full ML runtime
(TFLM, Caffe, or similar). The accuracy cost of the classical
approach, visible in Table~\ref{tab:overall}, is the explicit price
of removing those dependencies entirely.

The Edge model's 94~KB footprint warrants comment. A 10-tree random
forest with \texttt{max\_depth=8} over 24-dimensional input produces
substantially more total nodes than our earlier estimates suggested,
reflecting the interaction between tree depth, feature dimensionality,
and bootstrap-sample diversity at the given class balance. At 94~KB,
the Edge model requires a target with at least 128~KB flash — common
in STM32L4 and nRF52840 series devices but not the smallest Cortex-M0
parts. The Tiny model, at 3{,}656~bytes, remains compatible with
devices as small as 32~KB flash and represents a more aggressive
size--accuracy tradeoff: 0.708 accuracy at a flash cost three orders
of magnitude smaller than a typical TFLM-based neural network.

\subsection{Limitations}

Hardware deployment has not yet been validated; the C models have been
compiled and flash-profiled but not executed on physical silicon.
Additionally, the pipeline processes one beat at a time from pre-segmented
inputs; real-time deployment requires a preceding R-peak detector and
ring-buffer management, which are not included in the current pipeline.
Finally, the feature extraction step (particularly the FFT for
frequency-domain features) has not been benchmarked for latency on an
actual Cortex-M4 — only the inference function flash footprint is measured.

% ── VI. Conclusion ────────────────────────────────────────────────────────────
\section{Conclusion}
We described a reproducible pipeline for inter-patient ECG arrhythmia
classification targeting microcontroller deployment. The pipeline implements
the de~Chazal DS1/DS2 inter-patient split, extracts 24 features combining
morphological, spectral, and RR-interval information, trains three
classifier tiers spanning 121 to 283{,}328 tree nodes, exports the two
smaller tiers to verified dependency-free C, and measures their ARM
Cortex-M4 flash footprints from compiled object files. All code, trained
model artifacts, exported C files, and tests are available at
\url{https://github.com/bl4zeh3x/tinyml-explorations}.

The principal finding is that the N and V AAMI classes are well-handled
by classical feature-engineered classifiers under inter-patient evaluation
(F1 of 0.958 and 0.901 respectively for the Reference model), while S and
F remain difficult in a manner consistent with the broader literature.
Combining accurate V-class detection with a 3{,}656-byte inference
function (Tiny model) or a 94{,}252-byte function (Edge model) opens
practical paths to arrhythmia screening on wearable microcontroller
platforms of different flash-budget classes.

% ── References ────────────────────────────────────────────────────────────────
\bibliographystyle{IEEEtran}
\bibliography{refs}

\end{document}
